\section{Experimental Environment}
\label{sec:environment}

The environment models a layered enterprise network
(Fig.~\ref{fig:network}). An external attacker faces a gateway, service
hosts and a protected data store. The attacker wins by exfiltrating from
the store; the defender wins by detecting, containing and recovering first.

\begin{figure}[!t]
\centering
\begin{tikzpicture}[
  font=\scriptsize,
  host/.style={draw, rounded corners=1pt, minimum width=1.55cm, minimum height=0.55cm, align=center},
  arr/.style={-{Latex}}
]
\node[host] (inet) {Internet};
\node[host, below=0.35cm of inet] (gw) {Gateway};
\node[host, below=0.45cm of gw, xshift=-1.7cm] (sa) {Service A};
\node[host, below=0.45cm of gw] (sb) {Service B};
\node[host, below=0.45cm of gw, xshift=1.7cm] (sc) {Service C};
\node[host, below=0.55cm of sb, minimum width=5.1cm] (ds) {Data store};
\draw[arr] (inet) -- (gw);
\draw[arr] (gw) -- (sa);
\draw[arr] (gw) -- (sb);
\draw[arr] (gw) -- (sc);
\draw[arr] (sa) -- (ds);
\draw[arr] (sb) -- (ds);
\draw[arr] (sc) -- (ds);
\end{tikzpicture}
\caption{Layered cyber environment (Section~\ref{sec:environment}).
Attackers start outside the gateway under partial observability. The
terminal offensive objective is exfiltration from the data store.}
\label{fig:network}
\end{figure}


\paragraph{Tools and kill chain.}
Agents interact through a fixed catalogue of structured offensive and
defensive tools (reconnaissance through exfiltration; monitoring through
ingress control). Each tool has a resource cost and noise profile; parameters
are schema-validated before taking effect. Exfiltration requires composing a
multi-step chain under partial observability: reconnaissance, foothold,
optional credential access, lateral movement, then exfiltration. Transition
legality is enforced by the environment API, not by agent scripts. Tool
definitions and JSON schemas are maintained in the open-source artefact
(\hyperref[sec:code]{Code Availability}).

\paragraph{Partial observability and detection.}
Attackers begin with no host knowledge; defenders observe suspicion only on
monitored hosts. Noisy actions increase suspicion; detection fires with
probability $1-e^{-\lambda \sigma}$. Suspicion and monitoring intensity decay
each turn, which makes low-noise pauses after high-noise actions viable even
when objectives never mention waiting (Section~\ref{sec:results-ladder}).

\paragraph{Backends and defaults.}
All reported experiments use the in-memory simulator for reproducibility; an
optional Kubernetes backend implements the same interface. Default instances
use three service hosts, two vulnerabilities per host, twenty turns per
episode, and no inter-agent communication unless noted in
Table~\ref{tab:matrix}.
